\documentclass[12pt]{article}
\usepackage[utf8]{inputenc}
\usepackage[spanish]{babel}
\usepackage{amsmath}
\usepackage{amsfonts}
\usepackage{amssymb}
\usepackage{graphicx}
\usepackage{geometry}
\usepackage{tikz}
\usepackage{pgfplots}
\usepackage{booktabs}
\usepackage{array}
\usepackage{caption}
\usepackage{subcaption}
\usepackage{fancyhdr}
\usepackage{setspace}
\pgfplotsset{compat=1.18}
\geometry{margin=1in}
\setstretch{1.15}
\pagestyle{fancy}
\fancyhf{}
\fancyhead[L]{\nouppercase{\leftmark}}
\fancyhead[R]{Autocorrelación}
\fancyfoot[C]{\thepage}
\renewcommand{\headrulewidth}{0.4pt}
\renewcommand{\footrulewidth}{0pt}

\title{\textbf{Notas de Econometría} \\ \large Autocorrelación: Teoría, Detección y Corrección \\ \normalsize Basado en Wooldridge (2019) y Gujarati \& Porter (2010)}
\author{Clase: 17 de marzo de 2026}
\date{}

\begin{document}

\maketitle
\thispagestyle{empty}
\newpage
\setcounter{page}{1}
\tableofcontents
\newpage

\section{Introducción}
En el modelo clásico de regresión lineal, uno de los supuestos fundamentales de Gauss-Markov es la ausencia de correlación entre los términos de error para diferentes observaciones. Formalmente,
\begin{equation}
    \text{Cov}(u_i, u_j) = \mathbb{E}[u_i u_j] = 0 \quad \forall i \neq j.
\end{equation}
Cuando este supuesto se viola, se dice que existe \textbf{autocorrelación} (o correlación serial). La presencia de autocorrelación es particularmente frecuente en datos de series temporales, aunque también puede aparecer en datos de corte transversal (autocorrelación espacial) y en datos de panel.

El objetivo de estas notas es analizar las causas, consecuencias, métodos de detección y corrección de la autocorrelación, manteniendo un enfoque riguroso y aplicado.

\section{Naturaleza de la autocorrelación}
La autocorrelación puede modelarse mediante distintos procesos estocásticos. El más común es el proceso autorregresivo de orden uno, denotado AR(1):
\begin{equation}
    u_t = \rho u_{t-1} + \varepsilon_t, \quad |\rho| < 1,
\end{equation}
donde $\varepsilon_t$ es ruido blanco con media cero y varianza constante $\sigma^2_{\varepsilon}$. El parámetro $\rho$ mide la intensidad y dirección de la correlación:
\begin{itemize}
    \item $\rho > 0$: autocorrelación positiva (valores cercanos de $u_t$ tienden a tener el mismo signo).
    \item $\rho < 0$: autocorrelación negativa (alternancia de signos).
    \item $\rho = 0$: ausencia de autocorrelación.
\end{itemize}

Para procesos de orden superior, AR(p), la autocorrelación se extiende a rezagos mayores.

\section{Causas de la autocorrelación en los errores}
\subsection{Inercia o tendencia común}
En series de tiempo, las variables económicas suelen mostrar trayectorias suaves (PIB, inflación, desempleo). Esta inercia se transmite a los residuos si el modelo no captura completamente la dinámica subyacente.

\subsection{Sesgo de especificación: omisión de variables relevantes}
Cuando una variable relevante con comportamiento sistemático queda excluida del modelo, su efecto es absorbido por el término de error, generando autocorrelación. Por ejemplo, si se estima
\[
    I_t = \beta_1 + \beta_2 \text{Ingreso}_t + \beta_3 \text{Impuestos}_t + u_t
\]
pero la verdadera relación incluye la tasa de interés $r_t$, entonces $u_t$ contendrá $\alpha_4 r_t$ y, si $r_t$ está correlacionada en el tiempo, también lo estará $u_t$.

\subsection{Omisión de variables rezagadas (efectos dinámicos)}
Ignorar que el pasado influye en el presente produce autocorrelación. Un ejemplo clásico es la oferta agrícola:
\begin{align*}
    \text{Incorrecto:} & \quad S_t = \beta_1 + \beta_2 P_t - \beta_3 C_t + u_t \\
    \text{Correcto:} & \quad S_t = \beta_1 + \beta_2 P_{t-1} - \beta_3 C_{t-1} + u_t,
\end{align*}
donde $P_t$ es el precio y $C_t$ los costos de producción. La decisión de siembra depende de precios y costos pasados, no de los contemporáneos.

\subsection{Transformación de datos}
El uso de primeras diferencias para eliminar tendencias o para corregir colinealidad puede inducir autocorrelación en los errores transformados si el proceso original ya era ruido blanco.

\subsection{Manipulación de datos}
Procedimientos como interpolación o suavizamiento de datos faltantes introducen dependencia artificial entre observaciones consecutivas.

\subsection{Fenómeno de la telaraña (cobweb phenomenon)}
En mercados agrícolas, la oferta reacciona al precio del periodo anterior con un retraso, generando un ciclo de ajuste que se refleja en residuos correlacionados.

\section{Consecuencias de ignorar la autocorrelación}
Si se aplica MCO a pesar de existir autocorrelación, se producen los siguientes efectos:
\begin{enumerate}
    \item Los estimadores MCO siguen siendo insesgados y consistentes, pero dejan de ser eficientes (no son MELI).
    \item La matriz de varianzas-covarianzas estimada por MCO es sesgada, por lo que los errores estándar son incorrectos.
    \item Los contrastes $t$ y $F$ pierden validez, aumentando el riesgo de conclusiones erróneas.
    \item Las predicciones no son óptimas (mayor varianza).
\end{enumerate}

\section{Detección de la autocorrelación}
\subsection{Métodos gráficos}
Un primer paso es examinar los residuos $e_t$ frente al tiempo. La Figura \ref{fig:residuos_ar1} muestra un patrón típico de autocorrelación positiva generado a partir de un proceso AR(1) con $\rho = 0.8$. Se observa que los residuos se mantienen por encima o por debajo de cero durante períodos prolongados, evidenciando la persistencia.

\begin{figure}[h]
\centering
\begin{tikzpicture}
\begin{axis}[
    width=12cm,
    height=6cm,
    xlabel={$t$ (índice de tiempo)},
    ylabel={$e_t$ (residuos MCO)},
    axis lines=center,
    grid=major,
    grid style={dashed,gray!50},
    xtick={0,10,20,30,40,50},
    ytick={-2,-1,0,1,2},
    title={Residuos con autocorrelación positiva AR(1), $\rho = 0.8$},
    legend pos=north east,
]
\addplot+[black, thick, mark=none, smooth] coordinates {
    (0,0.2) (1,0.6) (2,1.0) (3,1.1) (4,1.2) (5,1.0) (6,0.8) (7,0.5) (8,0.3)
    (9,0.1) (10,-0.1) (11,-0.4) (12,-0.7) (13,-0.9) (14,-1.0) (15,-0.8) (16,-0.5)
    (17,-0.2) (18,0.0) (19,0.3) (20,0.5) (21,0.7) (22,0.9) (23,0.8) (24,0.6)
    (25,0.4) (26,0.1) (27,-0.2) (28,-0.5) (29,-0.7) (30,-0.9) (31,-0.6) (32,-0.3)
    (33,0.0) (34,0.2) (35,0.4) (36,0.6) (37,0.5) (38,0.3) (39,0.1) (40,-0.2)
    (41,-0.4) (42,-0.6) (43,-0.5) (44,-0.3) (45,-0.1) (46,0.1) (47,0.2) (48,0.4)
    (49,0.5) (50,0.3)
};
\addplot[black, dashed, thick, domain=0:50, samples=2] {0};
\legend{$e_t$, $e_t = 0$}
\end{axis}
\end{tikzpicture}
\caption{Residuos de un modelo estimado por MCO cuando el verdadero proceso de error es AR(1) con $\rho=0.8$. La línea punteada representa la media cero. La persistencia de los residuos por encima y por debajo de cero indica autocorrelación positiva.}
\label{fig:residuos_ar1}
\end{figure}

Otra herramienta gráfica fundamental es el correlograma (función de autocorrelación muestral, ACF). La Figura \ref{fig:acf} muestra las autocorrelaciones muestrales $r_k$ para los primeros 10 rezagos junto con bandas de confianza al 95\% bajo la hipótesis nula de ruido blanco. Las barras que exceden las bandas indican autocorrelación significativa en ese rezago. En este ejemplo, los rezagos 1 a 4 son significativos, y el decaimiento gradual es característico de un proceso AR(1).

\begin{figure}[h]
\centering
\begin{tikzpicture}
\begin{axis}[
    width=12cm,
    height=6cm,
    xlabel={Rezago $k$},
    ylabel={Autocorrelación muestral $r_k$},
    ybar,
    bar width=0.2cm,
    xtick={1,2,3,4,5,6,7,8,9,10},
    ytick={-0.5,0,0.5,1},
    ymin=-0.6, ymax=1,
    legend pos=north east,
    title={Correlograma de residuos (modelo AR(1) con $\rho=0.8$)},
    nodes near coords,
    nodes near coords style={anchor=south, yshift=4pt, font=\tiny},
    every node near coord/.append style={font=\tiny}
]
\addplot[black, fill=gray!50] coordinates {
    (1,0.85) (2,0.72) (3,0.58) (4,0.44) (5,0.31) (6,0.19) (7,0.10) (8,0.03) (9,-0.02) (10,-0.05)
};
\addplot[black, dashed, domain=1:10, samples=2] {0.2};
\addplot[black, dashed, domain=1:10, samples=2] {-0.2};
\legend{$r_k$, Bandas 95\%}
\end{axis}
\end{tikzpicture}
\caption{Función de autocorrelación muestral (ACF) de los residuos. Las bandas horizontales punteadas representan los límites de significación al 95\% bajo la hipótesis de ruido blanco ($\pm 1.96/\sqrt{T}$, con $T=50$). Las barras que sobresalen indican autocorrelación significativa. El decaimiento gradual es típico de un proceso AR(1).}
\label{fig:acf}
\end{figure}

\subsection{Contraste de Durbin-Watson}
El estadístico de Durbin-Watson se define como
\begin{equation}
    d = \frac{\sum_{t=2}^{T} (e_t - e_{t-1})^2}{\sum_{t=1}^{T} e_t^2},
\end{equation}
y aproximadamente $d \approx 2(1-\hat{\rho})$, donde $\hat{\rho}$ es el coeficiente de autocorrelación de primer orden estimado. La Figura \ref{fig:durbin} ilustra las regiones de decisión para un contraste bilateral al 5\% de significación, con valores críticos $d_L$ y $d_U$ que dependen del número de observaciones y de variables explicativas.

\begin{figure}[h]
\centering
\begin{tikzpicture}[scale=1]
\draw[->] (0,0) -- (10,0) node[right] {$d$};
\draw[thick] (0,0.2) -- (0,-0.2) node[below] {0};
\draw[thick] (2,0.2) -- (2,-0.2) node[below] {$d_L$};
\draw[thick] (3,0.2) -- (3,-0.2) node[below] {$d_U$};
\draw[thick] (6,0.2) -- (6,-0.2) node[below] {$4-d_U$};
\draw[thick] (7,0.2) -- (7,-0.2) node[below] {$4-d_L$};
\draw[thick] (10,0.2) -- (10,-0.2) node[below] {4};
\draw[<->] (0,0.5) -- (2,0.5) node[midway,above] {Autocorr. positiva};
\draw[<->] (2,0.5) -- (3,0.5) node[midway,above] {Indecisión};
\draw[<->] (3,0.5) -- (6,0.5) node[midway,above] {No autocorr.};
\draw[<->] (6,0.5) -- (7,0.5) node[midway,above] {Indecisión};
\draw[<->] (7,0.5) -- (10,0.5) node[midway,above] {Autocorr. negativa};
\node at (1.5, -0.8) {Rechazar $H_0$};
\node at (2.5, -0.8) {Región};
\node at (4.5, -0.8) {Aceptar $H_0$};
\node at (6.5, -0.8) {Región};
\node at (8.5, -0.8) {Rechazar $H_0$};
\end{tikzpicture}
\caption{Regiones de decisión para el estadístico de Durbin-Watson. $d_L$ y $d_U$ son los valores críticos inferior y superior que dependen del tamaño muestral $T$ y del número de regresores $k$ (sin incluir el término constante). Si $d < d_L$ se rechaza la hipótesis nula de no autocorrelación en favor de autocorrelación positiva; si $d > 4-d_L$ se rechaza en favor de autocorrelación negativa.}
\label{fig:durbin}
\end{figure}

\subsection{Contraste de Breusch-Godfrey (LM)}
A diferencia de Durbin-Watson, la prueba de Multiplicador de Lagrange (LM) permite detectar autocorrelación de orden superior ($AR(p)$) y es válida incluso si existen variables dependientes rezagadas ($Y_{t-1}$).

\textbf{Procedimiento:}
\begin{enumerate}
    \item Estimar el modelo original por MCO y obtener los residuos $\hat{u}_t$.
    \item Correr la regresión auxiliar:
    \begin{equation}
        \hat{u}_t = \alpha_1 + \alpha_2 X_t + \rho_1 \hat{u}_{t-1} + \rho_2 \hat{u}_{t-2} + \dots + \rho_p \hat{u}_{t-p} + \varepsilon_t
    \end{equation}
    \item Probar la hipótesis nula $H_0: \rho_1 = \rho_2 = \dots = \rho_p = 0$.
    \item El estadístico de prueba es $(n-p)R^2 \sim \chi^2_p$.
\end{enumerate}

\section{Medidas remediales: corrección de la autocorrelación}
Si se confirma la existencia de autocorrelación pura (no debida a sesgo de especificación), existen dos caminos principales:

\subsection{Mínimos cuadrados generalizados (MCG)}
Si conocemos la estructura de la autocorrelación (ej. un modelo $AR(1)$ con $\rho$ conocido), transformamos las variables para que el error sea ruido blanco.

\textbf{Transformación de Cochrane-Orcutt:}
Se define un nuevo modelo transformado:
\begin{equation}
    Y^*_t = \beta_1^* + \beta_2 X^*_t + \varepsilon_t
\end{equation}
Donde:
\begin{itemize}
    \item $Y^*_t = Y_t - \rho Y_{t-1}$
    \item $X^*_t = X_t - \rho X_{t-1}$
    \item $\beta_1^* = \beta_1(1-\rho)$
\end{itemize}
Este estimador es \textbf{MELI} (Mejor Estimador Lineal Insesgado).

\subsection{Errores estándar robustos (Newey-West)}
Si no queremos transformar el modelo o no conocemos la forma exacta de la autocorrelación, usamos el método de \textbf{HAC} (Heteroscedasticity and Autocorrelation Consistent).
\begin{itemize}
    \item Los coeficientes $\hat{\beta}$ siguen siendo los de MCO.
    \item Se corrigen las desviaciones estándar para que las pruebas $t$ y $F$ sean válidas.
    \item Es la opción preferida en software moderno (Stata, R, Python) para muestras grandes.
\end{itemize}

\subsection{Resumen comparativo de pruebas}
\begin{table}[h]
\centering
\caption{Comparación de pruebas de autocorrelación}
\label{tab:comparacion}
\begin{tabular}{lll}
\toprule
\textbf{Prueba} & \textbf{Ventajas} & \textbf{Limitaciones} \\
\midrule
Durbin-Watson & Fácil de calcular, muy común. & Solo AR(1), no admite $Y_{t-1}$. \\
Test de Rachas & No paramétrico (no requiere Normalidad). & Menos potente estadísticamente. \\
Breusch-Godfrey & Admite órdenes superiores y rezagos. & Requiere muestras grandes. \\
\bottomrule
\end{tabular}
\end{table}

\section{Ilustraciones gráficas adicionales}
A continuación se presentan tres gráficos detallados que complementan el análisis de la autocorrelación.

\subsection{Patrones de residuos: autocorrelación positiva vs. negativa}
La Figura \ref{fig:patrones} muestra dos patrones extremos: (a) autocorrelación positiva, donde los residuos consecutivos tienden a tener el mismo signo, generando rachas largas; (b) autocorrelación negativa, donde los residuos alternan signo sistemáticamente. Ambos casos violan el supuesto de independencia.

\begin{figure}[h]
\centering
\begin{subfigure}[b]{0.48\textwidth}
\centering
\begin{tikzpicture}
\begin{axis}[
    width=\textwidth,
    height=5cm,
    xlabel={$t$},
    ylabel={$e_t$},
    axis lines=center,
    grid=major,
    grid style={dashed,gray!50},
    ymin=-1.5, ymax=1.5,
    xmin=0, xmax=20,
    title={(a) Autocorrelación positiva}
]
\addplot[only marks, mark=*, mark size=1.5pt, black] coordinates {
    (1,0.4) (2,0.8) (3,1.1) (4,0.9) (5,0.6) (6,0.1) (7,-0.4) (8,-0.8) (9,-1.1) (10,-0.9)
    (11,-0.5) (12,0.1) (13,0.6) (14,1.0) (15,1.1) (16,0.7) (17,0.2) (18,-0.3) (19,-0.7) (20,-1.0)
};
\addplot[black, dashed, thick, domain=0:20, samples=2] {0};
\end{axis}
\end{tikzpicture}
\end{subfigure}
\hfill
\begin{subfigure}[b]{0.48\textwidth}
\centering
\begin{tikzpicture}
\begin{axis}[
    width=\textwidth,
    height=5cm,
    xlabel={$t$},
    ylabel={$e_t$},
    axis lines=center,
    grid=major,
    grid style={dashed,gray!50},
    ymin=-1.5, ymax=1.5,
    xmin=0, xmax=20,
    title={(b) Autocorrelación negativa}
]
\addplot[only marks, mark=*, mark size=1.5pt, black] coordinates {
    (1,0.8) (2,-0.7) (3,0.9) (4,-1.1) (5,0.6) (6,-0.9) (7,1.1) (8,-0.8) (9,0.7) (10,-1.2)
    (11,0.5) (12,-0.6) (13,0.9) (14,-1.0) (15,0.4) (16,-0.7) (17,1.0) (18,-0.8) (19,0.3) (20,-0.9)
};
\addplot[black, dashed, thick, domain=0:20, samples=2] {0};
\end{axis}
\end{tikzpicture}
\end{subfigure}
\caption{Patrones de residuos que indican autocorrelación: (a) rachas largas del mismo signo sugieren autocorrelación positiva; (b) alternancia sistemática de signos sugiere autocorrelación negativa. La línea punteada representa la media cero.}
\label{fig:patrones}
\end{figure}

\subsection{Sesgo de especificación: modelo lineal vs. verdadero cuadrático}
La Figura \ref{fig:especificacion} ilustra cómo la omisión de un término cuadrático genera residuos con un patrón sistemático. En el panel superior se muestra la nube de puntos de la relación verdadera (con forma de U) junto con la recta de regresión lineal estimada. El panel inferior muestra los residuos resultantes, que presentan una clara forma parabólica, violando el supuesto de media cero condicional y generando autocorrelación espuria.

\begin{figure}[h]
\centering
\begin{tikzpicture}
\begin{axis}[
    width=12cm,
    height=6cm,
    xlabel={$X$ (producción)},
    ylabel={$Y$ (costo marginal)},
    axis lines=left,
    grid=major,
    grid style={dashed,gray!50},
    ymin=0, ymax=10,
    xmin=0, xmax=10,
    title={Relación verdadera (curva en U) y ajuste lineal}
]
\addplot[only marks, mark=*, mark size=1.2pt, black!70] {0.4*(x-5)^2 + 2 + 0.5*rand};
\addplot[domain=1:9, thick, black] {0.5*x + 3.5};
\addplot[domain=1:9, dashed, thick, black] {0.4*(x-5)^2 + 2};
\legend{Datos observados, Regresión lineal (mal especificada), Relación verdadera}
\end{axis}
\end{tikzpicture}
\vspace{0.5cm}
\begin{tikzpicture}
\begin{axis}[
    width=12cm,
    height=4cm,
    xlabel={$X$},
    ylabel={Residuos $e$},
    axis lines=center,
    grid=major,
    grid style={dashed,gray!50},
    ymin=-2, ymax=2,
    xmin=0, xmax=10,
    title={Residuos del modelo lineal: patrón parabólico}
]
\addplot[only marks, mark=*, mark size=1.2pt, black!70] coordinates {
    (1, -0.5) (1.5, -1.2) (2, -1.6) (2.5, -1.9) (3, -1.8) (3.5, -1.5) (4, -1.0) (4.5, -0.3)
    (5, 0.5) (5.5, 1.3) (6, 1.7) (6.5, 1.9) (7, 1.6) (7.5, 1.0) (8, 0.4) (8.5, -0.2) (9, -0.9)
};
\addplot[domain=1:9, dashed, thick, black] {0.2*(x-5)^2 - 0.5};
\end{axis}
\end{tikzpicture}
\caption{Consecuencias de la omisión de una variable cuadrática. Superior: la verdadera relación es $Y=0.4(X-5)^2+2$ (curva en U), pero se ajusta un modelo lineal $Y=\beta_1+\beta_2X+u$. Inferior: los residuos muestran un patrón sistemático que indica mala especificación y puede inducir autocorrelación espuria.}
\label{fig:especificacion}
\end{figure}

\subsection{Comparación de intervalos de confianza: MCO vs. MCG}
La Figura \ref{fig:intervalos} muestra las distribuciones del estimador $\hat{\beta}_2$ bajo dos escenarios: (i) cuando se utiliza MCO ignorando autocorrelación, y (ii) cuando se aplica MCG con la transformación correcta. Bajo MCO, la varianza está sobreestimada, generando intervalos de confianza más amplios de lo necesario (menor precisión). En cambio, MCG produce estimadores eficientes (mínima varianza) y, por lo tanto, intervalos más estrechos, mejorando la potencia de las pruebas de hipótesis.

\begin{figure}[h]
\centering
\begin{tikzpicture}
\begin{axis}[
    width=12cm,
    height=6cm,
    axis lines=left,
    xtick={0},
    xticklabels={0},
    ytick=\empty,
    xmin=-5, xmax=6,
    ymin=0, ymax=1,
    xlabel={$\hat{\beta}_2$},
    ylabel={Densidad},
    title={Distribuciones de $\hat{\beta}_2$ bajo MCO (ignorando autocorrelación) y MCG (corrección)},
    legend pos=north west
]
\addplot[domain=-5:5, samples=100, thick, black] {exp(-x^2/8)/sqrt(8*pi)};
\addplot[domain=-5:5, samples=100, thick, dashed, black] {exp(-x^2/2)/sqrt(2*pi)};
\draw[<->, thick, black] (axis cs:-3.5,0.1) -- (axis cs:3.5,0.1) node[midway,below] {IC 95\% MCO};
\draw[<->, thick, dashed] (axis cs:-1.8,0.25) -- (axis cs:1.8,0.25) node[midway,above] {IC 95\% MCG};
\legend{MCO, MCG}
\end{axis}
\end{tikzpicture}
\caption{Comparación de la distribución de $\hat{\beta}_2$ cuando los errores siguen un proceso AR(1) con $\rho=0.8$. La curva continua (MCO) tiene mayor varianza, lo que conduce a intervalos de confianza más anchos. La curva punteada (MCG) es más eficiente (menor varianza) y sus intervalos de confianza son más estrechos, reflejando una estimación más precisa.}
\label{fig:intervalos}
\end{figure}

\section{Ejemplo ilustrativo}
Supongamos un modelo de demanda de dinero:
\[
    M_t = \beta_1 + \beta_2 \text{PIB}_t + \beta_3 r_t + u_t,
\]
con datos trimestrales. Al estimar por MCO se obtiene $d = 0.45$, indicando fuerte autocorrelación positiva. Se aplica el procedimiento de Cochrane-Orcutt, obteniendo $\hat{\rho}=0.78$ y coeficientes corregidos. Las pruebas de Breusch-Godfrey para rezagos 1 a 4 confirman la mejora.

\section{Conclusiones}
La autocorrelación es un problema frecuente en econometría que invalida los procedimientos de inferencia estándar. Su detección debe basarse en pruebas formales como Durbin-Watson o Breusch-Godfrey, complementadas con análisis gráfico. Las estrategias de corrección incluyen transformaciones GLS, métodos iterativos y el uso de errores estándar robustos. Una correcta especificación del modelo (incluyendo variables dinámicas) previene en gran medida su aparición.

\newpage

\begin{thebibliography}{9}
\bibitem{wooldridge} Wooldridge, J. M. (2019). \textit{Introducción a la econometría: un enfoque moderno}. Cengage Learning.
\bibitem{gujarati} Gujarati, D. N., \& Porter, D. C. (2010). \textit{Econometría}. McGraw-Hill.
\end{thebibliography}

\end{document}
